% Configurazione
\documentclass{article}

\usepackage{titling} % Required for inserting the subtitle
\usepackage{graphicx} % Reting images
\usepackage{tabularx} % Per l'ambiente tabularx (tabelle)
\usepackage{calc} % Sempre per le tabelle
\usepackage[hidelinks]{hyperref} % Per i collegamenti ipertestuali, ad esempio sulla table of contents
\usepackage[italian]{babel} % Per la lingua italiana nelle scritte automatiche
\usepackage{xcolor} % Per co
\usepackage{colortbl} % Per colorare le celle delle tabelle
\usepackage{lipsum} % Per generare lorem ipsum
\usepackage[normalem]{ulem} % Per sottotesto
\usepackage{array} % Per la vi fluttuante drisposte
\usepackage{ragged2e} % Pacchetto necessario per \justifying che giustifica il testo di tabelle

\newcommand{\ulhref}[2]{\href{#1}{\uline{#2}}} % Nuovo comando per sottolineare i link
\newcommand{\ulref}[1]{\uline{\ref{#1}}} % Nuovo comando per sottolineare i collegamenti a immagini
\setlength{\parindent}{0pt} % Rco dei paragrafi

\graphicspath{ {immagini/} {../../shared/immagini/} }


% Struttura
\begin{document}

\thispagestyle{plain} % Niente intestazione e piè di pagina


\begin{tikzpicture}[remember picture, overlay]
    % Punto di partenza al centro orizzontale nella metà superiore
    \coordinate (top_center) at ($(current page.north)!0.3!(current page.south)$);

    % UniPD: Logo e descrizione
    \node at (top_center) [anchor=north, xshift=-3cm, yshift=4.85cm] 
        {\includegraphics[width=0.15\textwidth]{Logo Universita di Padova.png}};
    \node at (top_center) [anchor=north, xshift=1.7cm, yshift=4.5cm]
        {\textcolor{red}{\textbf{Università degli Studi di Padova}}};
    \node at (top_center) [anchor=north, xshift=1.7cm, yshift=4cm]
        {\textcolor{red}{Laurea: Informatica}};
    \node at (top_center) [anchor=north, xshift=1.7cm, yshift=3.5cm]
        {\textcolor{red}{Corso: Ingegneria del Software}};
    \node at (top_center) [anchor=north, xshift=1.7cm, yshift=3cm]
        {\textcolor{red}{Anno Accademico: 2024/2025}};

    % SWEg Labs: Logo e descrizione
    \node at (top_center) [anchor=north, xshift=-2.85cm, yshift=1.5cm] 
        {\includegraphics[width=0.16\textwidth]{Logo SWEg.png}};
    \node at (top_center) [anchor=north, xshift=1.7cm, yshift=0.5cm]
        {\textbf{Gruppo: SWEg Labs}};
    \node at (top_center) [anchor=north, xshift=1.7cm, yshift=0cm]
        {Email: \textsf{gruppo.sweg@gmail.com}};
\end{tikzpicture}


\vspace{10cm}

{
\centering
\Huge\bfseries Glossario\par
\vspace{0.5cm}
\Large Versione 2.0.0\par
}

\vspace{2cm}
% Intestazione
\fancyhead[L]{1 \hspace{0.2cm} Informazioni generali} % Testo a sinistra

\pagenumbering{arabic} % Numerazione araba per il contenuto


\section{Informazioni generali}

\begin{itemize}
    \item \textbf{Tipo di riunione}: interna
    \item \textbf{Luogo}: meeting Discord
    \item \textbf{Data}: 06/02/2025
    \item \textbf{Ora inizio}: 17:30
    \item \textbf{Ora fine}: 19:00
    \item \textbf{Responsabile}: Giacomo Loat
    \item \textbf{Scriba}: Davide Verzotto
    \item \textbf{Partecipanti}:
    \begin{itemize}
        \item Federica Bolognini
        \item Michael Fantinato
        \item Giacomo Loat
        \item Filippo Righetto
        \item Riccardo Stefani
        \item Davide Verzotto
    \end{itemize}
\end{itemize}


\newpage
% Intestazione
\fancyhead[L]{Registro delle modifiche} % Testo a sinistra
\fancyhead[R]{\includegraphics[width=0.16\textwidth]{sweg_logo_sito_inverted.png}} % Immagine a destra

% Piè di pagina
\fancyfoot[L]{Piano di Qualifica}       % Testo a sinistra
\fancyfoot[C]{\thepage}                % Numero di pagina al centro
\fancyfoot[R]{Versione 1.0.0}          % Testo a destra

\pagenumbering{roman} % Numerazione romana per l'indice


\section*{Registro delle modifiche}

\begin{table}[h]
    \centering
    \rowcolors{2}{lightgray}{white}
    \begin{tabular}{|c|c|p{5cm}|p{3cm}|p{3cm}|}
        \hline
        \rowcolor[gray]{0.75}
        \textbf{Versione} & \textbf{Data} & \multicolumn{1}{|c|}{\textbf{Descrizione}} & 
        \multicolumn{1}{|c|}{\textbf{Autore}} & \multicolumn{1}{|c|}{\textbf{Verifica}}\\
        \hline
        1.0.0 & 23-01-25 & Approvazione del documento & Michael Fantinato & Michael Fantinato\\
        \hline
        0.3.0 & 23-01-25 & Verifica del documento & Davide Verzotto & Davide Verzotto\\
        \hline
        0.2.7 & 17-01-25 & Stesura paragrafo \S\bulref{subsec:Indice di Gulpease} & Giacomo Loat & Riccardo Stefani\\
        \hline
        0.2.6 & 17-01-25 & Stesura di gran parte della sezione \S\bulref{sec:Cruscotto di valutazione della qualità} & Riccardo Stefani & Giacomo Loat\\
        \hline
        0.2.5 & 17-01-25 & Stesura sezione \S\bulref{sec:Valutazioni per il miglioramento}
        e stesura paragrafo \S\bulref{subsec:Requisiti soddisfatti} & Giacomo Loat & Riccardo Stefani\\
        \hline
        0.2.4 & 15-01-25 & Stesura sezione \S\bulref{sec:Test di sistema} & Davide Verzotto & Riccardo Stefani\\
        \hline
        0.2.3 & 14-01-25 & Stesura sezione \S\bulref{sec:checklist} & Federica Bolognini & Davide Verzotto\\
        \hline
        0.2.2 & 12-01-25 & Stesura sezione \S\bulref{sec:Test di accettazione} & Filippo Righetto & Davide Verzotto\\
        \hline
        0.2.1 & 31-12-24 & Continuata stesura sezione \S\bulref{sec:Qualità di prodotto} & Davide Verzotto & Riccardo Stefani\\
        \hline
        0.2.0 & 06-12-24 & Verifica del documento allo stato attuale & Davide Verzotto & Davide Verzotto\\
        \hline
        0.1.2 & 02-12-24 & Continuata stesura sezione \S\bulref{sec:Piano di qualità} & Michael Fantinato & Riccardo Stefani\\
        \hline
        0.1.2 & 01-12-24 & Inizio stesura sezione \S\bulref{sec:Cruscotto di valutazione della qualità} & Riccardo Stefani & Giacomo Loat\\
        \hline
        0.1.1 & 09-11-24 & Inizio stesura sezioni \S\bulref{sec:introduzione}, \S\bulref{sec:Piano di qualità} e \S\bulref{sec:strategie di testing} & Michael Fantinato & Federica Bolognini\\
        \hline
        0.1.0 & 05-11-24 & Creazione del documento & Riccardo Stefani & Davide Verzotto\\
        \hline
    \end{tabular}
    \caption{Registro delle modifiche}
\end{table}

\newpage
\tableofcontents
\newpage

% Intestazione
\fancyhead[L]{1 \hspace{0.2cm} Introduzione} % Testo a sinistra

\pagenumbering{arabic} % Numerazione araba per il contenuto



\section{Introduzione}
\label{sec:introduzione}

\subsection{Scopo del documento}
Per i progetti di sviluppo software che puntano a soddisfare gli elevati 
standard qualitativi delineati nei principi dell'ingegneria del software, 
un documento dettagliato del piano di qualità è indispensabile.
Comprendere e valutare la \emph{qualità}\textsubscript{\textit{\textbf{G}}} 
del prodotto sono concetti fondamentali 
per effettuare confronti e determinare quanto un prodotto aderisce alle aspettative.
Questo documento mira a fornire una descrizione completa e precisa delle 
\emph{metriche}\textsubscript{\textit{\textbf{G}}} e delle 
metodologie utilizzate per il controllo e la misurazione 
della qualità nelle diverse componenti del software.
Definiremo gli obiettivi di qualità, i processi e le risorse necessarie per 
raggiungerli, oltre a specificare i test previsti con la relativa documentazione, 
comprese metodologie ed esiti.
Il documento sarà una risorsa preziosa per chi sviluppa il prodotto, per gli 
utenti finali e per coloro che ne valuteranno la qualità.

\subsection{Glossario}
Al fine di evitare eventuali equivoci o incomprensioni riguardo la terminologia 
utilizzata all’interno di questo documento, 
abbiamo valutato di adottare un Glossario, con file apposito, in cui
vengono riportate tutte le definizioni rigogliose delle parole ambigue utilizzate 
in ambito di
questo progetto.\\
Nel documento appena descritto verranno riportati tutti i termini definiti nel
loro ambiente di utilizzo con annessa descrizione del loro significato.\\
La presenza di un termine all'interno del Glossario sarà indicata con una "G" 
posizionata al pedice della parola.

\subsection{Maturità e miglioramenti}
Questo documento è stato redatto seguendo un approccio incrementale, 
con l'obiettivo di facilitare l'adattamento alle esigenze mutevoli, stabilite 
di comune accordo tra i membri del gruppo di progetto e l'azienda proponente.\\
Pertanto, il documento non può essere considerato definitivo o esaustivo, ma 
piuttosto un punto di partenza per un continuo aggiornamento e affinamento.

\subsection{Riferimenti}
\subsubsection{Riferimenti normativi}
\begin{itemize}
    \item \bulhref{https://sweg-labs.github.io/Documentazione/output/RTB/Documentazione\%20interna/norme_progetto_v1.0.0.pdf}{Norme di Progetto v.1.0.0}; \\
    \item \bulhref{https://www.math.unipd.it/~tullio/IS-1/2024/Progetto/C9.pdf}{Capitolato C9: BuddyBot}; \\
    \item \bulhref{https://www.math.unipd.it/~tullio/IS-1/2024/Dispense/PD1.pdf}{Regolamento progetto didattico}; \\
    \item \bulhref{https://www.iso.org/obp/ui\#iso:std:iso:31000:ed-2:v1:en}{Standard ISO/IEC 31000:2018}.
\end{itemize}

\subsubsection{Riferimenti informativi}
\begin{itemize}
    \item \bulhref{https://sweg-labs.github.io/Documentazione/output/RTB/Documentazione\%20esterna/glossario_v1.0.0.pdf}{Glossario v.1.0.0}; \\
    \item \bulhref{https://www.math.unipd.it/~tullio/IS-1/2024/Dispense/T7.pdf}{T7 - Qualità del software}; \\
    \item \bulhref{https://www.math.unipd.it/~tullio/IS-1/2024/Dispense/T8.pdf}{T8 - Qualità di processo}; \\
    \item \bulhref{https://www.math.unipd.it/~tullio/IS-1/2024/Dispense/T9.pdf}{T9 - Verifica e validazione: introduzione}; \\
    \item \bulhref{https://www.math.unipd.it/~tullio/IS-1/2024/Dispense/T10.pdf}{T10 - Verifica e validazione: analisti statica}; \\
    \item \bulhref{https://www.math.unipd.it/~tullio/IS-1/2024/Dispense/T11.pdf}{T11 - Verifica e validazione: analisti dinamica}. \\


\end{itemize}
\section{Valutazione del Capitolato selezionato}

\subsection{C9: BuddyBot}

\subsubsection{Descrizione}
Il capitolato richiede di sviluppare un chatBot capace di rispondere
a diversi tipi di richieste relative ai dati interni all'azienda, 
rendendo inoltre possibile per gli utenti ricevere informazioni riguardanti 
i contenuti presenti su piattaforme esterne.

\subsubsection{Dominio}
\paragraph{Dominio tecnologico:}
Viene richiesto l'utilizzo di un LLM open-source o di API di modelli 
a pagamento (come ChatGPT). Il modello deve poter comunicare tramite le API di Jira, 
Confluence e GitHub. Una volta recuperate le informazioni necessarie per la risposta,
deve elaborarle e formulare una risposta finale. Tutte le chat devono essere
salvate su un database e devono poter essere recuperate. L'azienda non richiede
l'uso di una tecnologia specifica, ma fornisce alcuni suggerimenti:
\begin{itemize}
    \item \textit{\textbf{OpenAI}}: Utilizzo tramite API per le funzionalità NLP, ovvero di comprensione e generazione del testo.
    \item \textit{\textbf{LangChain}}: Progetto open-source che permette l'integrazione di modelli AI senza richiedere una conoscenza dettagliata del loro funzionamento interno.
    \item \textit{\textbf{Angular}}: Framework front-end per la costruzione di applicazioni web.
    \item \textit{\textbf{Node/NestJS}}: Framework per lo sviluppo di applicazioni server-side orientato ai microservizi e alle API RESTful.
    \item \textit{\textbf{Spring Boot}}: Framework Java che offre strumenti per la creazione di applicazioni standalone, con supporto integrato per database, sicurezza e gestione delle dipendenze.
\end{itemize}


\paragraph{Dominio applicativo:}
BuddyBot vuole diventare un assistente a 360 gradi per qualsiasi membro 
dell'azienda. Si propone di essere facilmente accessibile tramite
un'interfaccia web, dove l'utente può formulare domande e recuperare lo storico delle vecchie chat.
Il bot deve poter rispondere a domande riguardanti file presenti in Jira, GitHub e Confluence, 
riuscendo anche a incrociare dati provenienti da diverse fonti. Il compito principale di 
BuddyBot sarà quindi aiutare i membri dell'azienda a trovare e combinare i dati in 
modo più rapido e preciso, fornendo risposte in linguaggio naturale.

\subsubsection{Motivazione della scelta}
La scelta finale è ricaduta su questo capitolato per diversi motivi:
\begin{itemize}
    \item \textit{\textbf{Disponibilità dell'azienda}}
    : L'azienda si è dimostrata da subito molto disponibile; le risposte alle 
    mail sono arrivate in brevissimo tempo e si è subito manifestata la volontà 
    di organizzare un incontro. Inoltre, si sono proposti di fornire anche 
    supporto tecnico durante lo sviluppo e degli spazi sulle loro piattaforme 
    dedicati a noi, qualora ce ne fosse bisogno.
    \item \textit{\textbf{Interesse}}
    : Il prodotto finale ha suscitato grande interesse nel gruppo, 
    soprattutto per le tecnologie con cui prevediamo di lavorare per 
    realizzarlo. In particolare, ci è sembrato il capitolato con il miglior 
    utilizzo dell'IA.
    \item \textit{\textbf{Fattibilità del progetto}}
    : Confrontando le richieste tra i vari capitolati, BuddyBot ci è 
    sembrato il miglior compromesso tra difficoltà e interesse.
    \item \textit{\textbf{Libertà tecnologica}}
    : La libertà lasciata nella scelta delle tecnologie ci è sembrata
     molto vantaggiosa, permettendoci di utilizzare tecnologie con le
      quali alcuni membri del gruppo hanno già esperienza.
\end{itemize}

\subsubsection{Conclusioni}
Visti i motivi sopra citati e non trovando in nessun altro capitolato una combinazione equivalente, 
il gruppo ha deciso di scegliere BuddyBot come progetto per cui candidarsi.

\section{Valutazione dei Capitolati preferiti}


\subsection{C7: LLM: Assistente digitale}

\subsubsection{Descrizione}
Il Progetto prevede la realizzazione di un assistente virtuale 
che aziende il cui core business è dato dalla vendita di prodotti
possono mettere a disposizione dei propri clienti, 
al fine di facilitare la ricerca di informazioni sui prodotti disponibili 
e rispondere alle domande più frequenti.
\subsubsection{Dominio}
\paragraph{Dominio tecnologico:}
Non ci sono tecnologie obbligatorie, ma sono state consigliate:
\begin{itemize}
    \item \textit{\textbf{MySQL}}: per la gestione del database relazionale.
    \item \textit{\textbf{LLM}}: il proponente ha indicato un insieme di modelli open source, tra questi si evidenzia Italia by iGenius, modello con 9 miliardi di parametri addestrato con un dataset al 90\% italiano.
    \item \textit{\textbf{API Rest}}: per la comunicazione tra il modello LLM e l'applicativo di interazione con l'utente.
    \item \textit{\textbf{ODBC}}: Open Database Connectivity, standard per la comunicazione da e per il database.
    \item \textit{\textbf{.NET MAUI}}: framework per lo sviluppo di applicazioni cross-platform.
\end{itemize}
\paragraph{Dominio applicativo:}
Il progetto si inserisce nel contesto della digitalizzazione e dell’uso del Machine Learning 
per analizzare grandi quantità di dati aziendali e sviluppare sistemi di interazione uomo-macchina avanzati. 
Nello specifico è rivolto alle aziende del settore della vendita di alimenti, 
dove la conoscenza dettagliata dei prodotti è spesso affidata a specialisti. 
Mira a creare un assistente virtuale che aiuti i clienti a trovare informazioni sui prodotti disponibili 
e risponda alle domande frequenti, migliorando l’accessibilità delle informazioni 
e l’efficienza del processo di assistenza clienti, riducendo la dipendenza dagli specialisti umani 
e offrendo risposte rapide e accurate alle domande dei clienti.
\subsubsection{Criticità riscontrate}
Fatta eccezione per l’elevata complessità delle tecnologie necessarie per lo sviluppo, 
che richiederà un certo periodo di studio da parte di tutti i membri del gruppo, 
non abbiamo riscontrato importanti criticità. 
Tuttavia ci sono alcuni aspetti che ci hanno convinto a scegliere il capitolato C9 piuttosto di questo. 
I due capitolati trattano progetti molto simili, ma riteniamo che per quanto riguarda il capitolato C9 
l’azienda proponente offra un supporto più adeguato alle nostre esigenze, 
e le tecnologie suggerite hanno suscitato maggiore interesse da parte del gruppo.
\subsubsection{Conclusioni}
Il gruppo ha apprezzato molto il progetto presentato, 
riconoscendone il valore e la potenziale utilità. 
Tuttavia, abbiamo deciso di non procedere con questo progetto, 
non per una mancanza di qualità, ma perché un altro capitolato ha suscitato un maggiore interesse da parte nostra.

\subsection{C2: Vimar GENIALE}

\subsubsection{Descrizione}
Nel capitolato si propone la realizzazione di un chatBot che gli installatori 
possono interpellare per reperire  informazioni testuali e grafiche sui prodotti 
Vimar presenti all’interno del sito ufficiale.
\subsubsection{Dominio}
\paragraph{Dominio tecnologico:}
Per lo svolgimento del capitolato ci sono delle tecnologie obbligatorie da utilizzare:
\begin{itemize}
    \item \textit{\textbf{Docker}}: Permette lo sviluppo di istanze isolate su cui testare il prodotto in un ambiente controllato
    \item \textit{\textbf{Git}}: Sistema di versionamento
    \item \textit{\textbf{Approcio RAG (retrieval Augmented Generation)}}: Permette il recupero di informazioni e generazione di risposte in linguaggio naturale
\end{itemize}
Oltre alle tecnologie obbligatorie abbiamo anche delle tecnologie consigliate:
\begin{itemize}
    \item \textit{\textbf{Angular}}:  Per realizzare il front-end della web app
    \item \textit{\textbf{Python}}:  Per realizzare le API e la business-logic
    \item \textit{\textbf{Scrapy}}: Utile per raccogliere dati da diversi siti web
    \item \textit{\textbf{PostgreSQL}}: Per database relazionale in cui immagazzinare i dati 
\end{itemize}
\paragraph{Dominio applicativo:}
Questo capitolato si propone per aiutare e facilitare il lavoro degli installatori 
con un applicativo web dove potranno contattare e reperire informazioni molto utili 
per agevolare l’installazione degli apparati.
Lo scopo di questo progetto è quello di realizzare un applicativo web simile come 
funzionamento ad un chat bot attraverso il quale gli installatori potranno chiedere 
in linguaggio naturale ed ottenere informazioni molto importanti come ad esempio: 
dati tecnici, schema elettrico, istruzioni funzionamento ecc.
  
\subsubsection{Criticità riscontrate}
\begin{itemize}
    \item \textit{\textbf{Minore libertà tecnologica}}:  L’azienda è più rigida nelle tecnologie che richiede di utilizzare nel capitolato rispetto ad altri proponenti.
    \item \textit{\textbf{Poca reperibilità}}:  L’azienda si è resa meno disponibile rispetto agli altri proponenti poiché non ha accettato di tenere un incontro telematico, preferendo un documento condiviso. Il che ci ha dato un segnale negativo riguardo al supporto che potremmo ricevere durante la fase progettuale.
\end{itemize}
\subsubsection{Conclusioni}
Il capitolato C2 è un’ottima occasione per chi vuole imparare competenze avanzate 
e arricchire il proprio profilo. Tra i vantaggi principali ci sono l’opportunità 
di fare esperienza sia nel front-end che nel back-end (full-stack) e l’uso di 
tecnologie moderne e molto richieste. Tuttavia, viste anche le criticità sopra citate, 
abbiamo deciso di non procedere con questo progetto perchè un altro capitolato ha 
suscitato un maggiore interesse da parte nostra.
\section{Valutazione degli altri Capitolati}


\subsection{C1: Artificial QI}

\subsubsection{Pro}

\begin{itemize}
    \item \textbf{Innovazione e utilità}: il gruppo trova interessante cercare di trovare una soluzione per il problema molto attuale della valutazione delle abilità degli LLM.
    \item \textbf{Punti di partenza}: il proponente Zucchetti è disponibile a fornire un programma esterno che espone delle API da testare, che è un'opportunità per avere una parte del sistema già pronta.
    \item \textbf{Molti requisiti opzionali}: sono consigliati vari spunti per estendere l'applicazione nel caso in cui si riesca a concludere in tempi brevi i requisiti obbligatori.
\end{itemize}

\subsubsection{Contro}

\begin{itemize}
    \item \textbf{Complessità del problema}: il capitolato tratta un problema aperto nel mondo dell'informatica, il che costituisce una complessità di ricerca esplorativa che va oltre lo sviluppo dell'applicativo e che non è calcolabile a priori.
    \item \textbf{Genericità degli strumenti}: vengono presentate molte poche tecnologie consigliate rispetto ad altri capitolati, e i suggerimenti rimangono molto generici a riguardo.
    \item \textbf{Dominio incognito e dataset assente}: non è chiaro quale sia il dominio di conoscenza che è da testare nell'LLM di turno, infatti non viene mai specificato che tipo di domande e risposte devono essere presenti nel dataset. Sembra inoltre che sia compito del gruppo scrivere le domande e risposte che facciano da training set, il che costituisce un onere non da poco considerata l'enorme quantità di dati necessari per allenare il sistema.
\end{itemize}


\subsection{C3:	Automatizzare le routine digitali tramite l’intelligenza generativa}

\subsubsection{Pro}

\begin{itemize}
    \item \textbf{Automazione delle Routine}: È interessante che il progetto punti ad automatizzare compiti ripetitivi come l’organizzazione delle e-mail e la gestione dei calendari, migliorando così l’efficienza e l'organizzazione aziendale.
    \item \textbf{Estendibilità}: L’applicativo è concepito in modo modulare, permettendo di aggiungere facilmente nuove automazioni e blocchi funzionali.
    \item \textbf{Supporto e Formazione}: Var Group fornisce supporto continuo per le tecnologie utilizzate, formazione e incontri periodici per garantire la qualità del progetto e per gestire l'evoluzione del software.
\end{itemize}

\subsubsection{Contro}

\begin{itemize}
    \item \textbf{Incremento dei tempi di sviluppo}: Potremmo dover affrontare ritardi dovuti alla necessità di studiare nuove e diverse tecnologie tra cui Python, Swift, NodeJS, MongoDB e Typescript. Questo potrebbe impattare la consegna nei tempi previsti, con conseguente aumento dei costi di progetto.
    \item \textbf{Scarso interesse}: il gruppo ha preferito concentrarsi su altri progetti con focus più accattivanti.
\end{itemize}


\subsection{C4:	NearYou - Smart custom advertising platform}

\subsubsection{Pro}

\begin{itemize}
    \item \textbf{Innovatività}: il progetto utilizza tecnologie avanzate come intelligenza artificiale generativa, data stream processing e targeting geografico in tempo reale. Questo lo rende un progetto ideale per acquisire esperienza con tecnologie molto richieste nel mercato attuale.
    \item \textbf{Flessibilità tecnologica}: il progetto lascia ampia libertà di scelta sugli strumenti e tecnologie da adottare.
    \item \textbf{Licenza}: la licenza del software sarà interamente nostra, potendo esporre il progetto nel curriculum.
\end{itemize}

\subsubsection{Contro}

\begin{itemize}
    \item \textbf{Complessità elevata}: il progetto richiede la gestione di argomenti forse troppo complessi data la scarsa esperienza dei membri del gruppo negli ambiti richiesti.
    \item \textbf{Rischio di dispersione}: la possibilità di scegliere le tecnologie da utilizzare potrebbe essere un'arma a doppio taglio, portando a dispersione o a scelte tecniche inadeguate, soprattutto se non si è esperti in tutti gli ambiti.
    \item \textbf{Scarso interesse}: il gruppo ha preferito concentrarsi su altri progetti con focus più accattivanti.
\end{itemize}


\subsection{C5:	3Dataviz}

\subsubsection{Pro}

\begin{itemize}
    \item \textbf{Gestione dei dati}: Il progetto è molto utile per quanto riguarda l’apprendimento, la gestione e la visualizzazione di un grande numero di dati in maniera chiara e interattiva.
    \item \textbf{Ampia libertà progettuale}: Il progetto offre grande libertà dal punto di vista della realizzazione dei vari grafici e delle funzionalità richieste. Questa flessibilità consente di personalizzare il sito in base alle esigenze e alle competenze del team.
\end{itemize}

\subsubsection{Contro}

\begin{itemize}
    \item \textbf{Richieste competenze avanzate}: La creazione e gestione di un sistema di visualizzazione 3D richiede delle competenze tecniche molto avanzate nel campo della grafica tridimensionale.
    \item \textbf{Studio di tecnologie nuove}: Il progetto richiede lo studio di tecnologie mai utilizzate prima, rendendo necessario un approfondimento e uno studio individuale per acquisire le competenze necessarie.
    \item \textbf{Scarso interesse}: il gruppo ha preferito concentrarsi su altri progetti con focus più accattivanti.
\end{itemize}


\subsection{C6: Sistema di gestione di un magazzino distribuito}

\subsubsection{Pro}

\begin{itemize}
    \item \textbf{Machine Learning}: Si richiede l'utilizzo del machine Learning per prevedere la richiesta futura di un certo materiale, materia che può essere molto interessante su cui lavorare.
    \item \textbf{Utilità}: Il problema della gestione del magazzino c’e’ da molti anni e colpisce molte aziende, imparare a creare software innovativi per la gestione di esso può essere molto utile e formativo.
\end{itemize}

\subsubsection{Contro}

\begin{itemize}
    \item \textbf{Complessità dei microservizi}: L'architettura basata su microservizi introduce maggiore complessità nella gestione del sistema e nell'orchestrazione, richiedendo più competenze e lavoro..
    \item \textbf{Sicurezza}: Richiesta di sicurezza avanzata tra il sistema centrale e i sistemi locali, non richiesto da altri capitolati. Questa richiesta aumenta la complessità del sistema, rendendo anche più difficile quantificare le ore di allenamento necessarie.
    \item \textbf{Requisiti di test e validazione}: La necessità di effettuare test estensivi e validazioni di algoritmi predittivi e di sicurezza, implica un elevato sforzo continuo.
\end{itemize}


\subsection{C8: Requirement Tracker - Plug-in VS Code}

\subsubsection{Pro}

\begin{itemize}
    \item \textbf{Supervisione continua}: disponibilità dell’azienda ad essere presente sia in presenza che in remoto, potendo così assistere il gruppo e guidarlo nella realizzazione.
    \item \textbf{Modelli di AI pre-addestrati}: è possibile utilizzare modelli di AI già esistenti e sfruttarli tramite API. Il team potrebbe quindi iniziare a lavorare senza dover programmare da zero un modello.
\end{itemize}

\subsubsection{Contro}

\begin{itemize}
    \item \textbf{Necessario studio di tecnologie sconosciute}: è necessario implementare un’architettura di tipo modulare, che il gruppo non ha ancora usato approfonditamente, dunque nonostante l’occasione, il suo apprendimento rallenterebbe il lavoro per il tempo necessario.
    \item \textbf{Vincoli stringenti}: i vincoli obbligatori imposti sono in quantità tale da presentare un ridotto spazio di manovra o di inventiva per il team.
    \item \textbf{Proof of Concept}: la necessità di sviluppare un P.o.C. già a metà progetto, può richiedere uno sforzo elevato, soprattutto essendo previste numerose tecnologie nuove da apprendere.
\end{itemize}

\end{document}
